\documentclass[11pt]{article}

\usepackage[utf8]{inputenc}
\usepackage[T1]{fontenc}
\usepackage{lmodern}
\usepackage[margin=1in]{geometry}
\usepackage{amsmath, amssymb}
\usepackage{microtype}
\usepackage{enumitem}
\usepackage{xcolor}
\usepackage{hyperref}
\hypersetup{colorlinks=true, linkcolor=black, urlcolor=blue}

% Number this section as §3 to match the larger paper.
\setcounter{section}{2}

% Convenience macros.
\newcommand{\anchorW}{W_{\text{anchor}}}
\newcommand{\anchorB}{b_{\text{anchor}}}
\newcommand{\Lewc}{L_{\text{EWC}}}
\newcommand{\Ltask}{L_{\text{task}}}
\newcommand{\reff}{r_{\text{eff}}}
\newcommand{\ac}{\text{ac}}
\newcommand{\R}{\mathbb{R}}

\title{Trioron --- Methods (\S3)}
\author{}
\date{}

\begin{document}
\maketitle

\section{Methods}
\label{sec:methods}

\paragraph{Notation.}
$n$ is the number of nodes in a layer; $d$ is the layer's incoming
dimension; $B$ is batch size; $t$ indexes tasks in a curriculum.

\subsection{The Trioron Node}
\label{sec:trioron-node}

A \emph{trioron} is a single computational unit characterized by three
coupled state variables, generalizing the standard artificial neuron:

\begin{itemize}[leftmargin=1.5em]
  \item $w \in \R^{d}$ --- incoming weights (a row of the layer's weight
        matrix).
  \item $\lambda \in \R_{\geq 0}$ --- per-node \emph{plasticity coefficient}:
        a per-unit scalar that scales the unit's contribution to the
        elastic-weight consolidation (EWC) penalty. Larger
        $\lambda \Rightarrow$ stiffer; smaller $\lambda \Rightarrow$ more
        plastic.
  \item $u \in \R$ --- per-node \emph{utility score}: a running estimate
        of the unit's contribution to the network's outputs. Used as the
        signal for pruning decisions (cells with persistently low $u$ are
        candidates for removal).
\end{itemize}

A \emph{trioron layer} aggregates $n$ trioron nodes that share an incoming
dimension. We hold the per-node state in vectors $\lambda, u \in \R^{n}$
and a weight matrix $W \in \R^{n \times d}$ where row $i$ is $w_i$. The
forward pass for input batch $X \in \R^{B \times d}$ is
\begin{equation}
  H = \sigma\!\bigl(X\,(r \odot W)^{\top} + b\bigr),
\end{equation}
where $\sigma$ is the activation function, $b \in \R^{n}$ is the bias, and
$r \in [0, 1]^{n}$ is the layer's \emph{routing scale} --- a per-node
multiplicative gain on incoming weights, defaulting to all ones. The
routing scale is the substrate for routing-starvation consolidation
(\S\ref{sec:dreaming-phase}).

Each layer additionally maintains EWC state: anchor weights $\anchorW$ and
$\anchorB$ snapshotted at task boundaries, and Fisher-information
accumulators $F_W, F_b$ updated as a running EMA of squared gradients.
Per-node $\lambda$ is derived from the row-mean of Fisher:
$\lambda_i = \operatorname{mean}_j F_{W,i,j}$.

\paragraph{Epigenetic analogue.}
$w$ is the synaptic-strength channel (modulated by experience-dependent
LTP/LTD); $\lambda$ is a per-cell modifier on plasticity, with biological
correlates including DNA-methylation status of plasticity-related genes
(e.g., BDNF) and perineuronal-net maturation, both of which gate how
readily a cell's synapses can be modified; $u$ is the cell-importance
signal that biological systems approximate via activity-dependent
neurotrophin release. The triplet $(w, \lambda, u)$ captures \emph{which
connections are present, how rigid they are, and how much the cell
matters} --- three quantities that biological systems modulate
independently.

\subsection{Continual-Learning Substrate (EWC with Per-Node Plasticity)}
\label{sec:ewc-substrate}

Following Kirkpatrick et~al. (2017), we apply elastic weight consolidation
to protect prior-task knowledge. The penalty for a single trioron layer is
\begin{equation}
  \Lewc = \sum_i \lambda_i \left[
    \sum_j (W_{ij} - W_{\text{anchor},ij})^{2}
         + (b_i - b_{\text{anchor},i})^{2}
  \right].
\end{equation}
The total loss is $L = \Ltask + \beta \cdot \Lewc$ for an EWC strength
$\beta$ set per-curriculum. We refresh $\lambda_i$ from Fisher at the end
of each task (after a batched re-estimation pass over recent task data),
then snapshot $\anchorW \leftarrow W$ as the new anchor. This separation
--- estimate Fisher at task end, then anchor --- follows the original EWC
recipe and avoids the EMA-drift artifacts of online updating during
training.

For longer curricula we optionally accumulate Fisher across tasks via the
Online EWC variant (Schwarz et~al. 2018):
\begin{equation}
  F_t \leftarrow \gamma \cdot F_{t-1} + F_{\text{new}},
  \qquad \gamma \in [0, 1],
\end{equation}
with $\gamma = 1$ reducing to vanilla EWC. We use Online EWC
$\gamma = 0.95$ as a competitor baseline in \S4.

A key modification: when an apoptosis spike fires
(\S\ref{sec:dreaming-phase}), each node's \emph{effective} stiffness is
scaled by $\max(0, 1 - p_i)$, where $p_i$ is the node's apoptosis pulse:
\begin{equation}
  \lambda_i^{\text{eff}} = \lambda_i \cdot \max(0, 1 - p_i).
\end{equation}
This temporarily reduces a surviving cell's EWC pinning when a neighbor
has just died, allowing it to absorb the dead neighbor's role.
Mechanistically analogous to acute glial-derived neurotrophic factor
release post-apoptosis, which transiently increases plasticity in adjacent
cells.

\subsection{Structural Plasticity}
\label{sec:structural-plasticity}

The architecture supports in-place growth and pruning of nodes during
training, with cross-layer consistency.

\paragraph{Cellular division (\texttt{grow\_node}).}
Adds one row to a layer's weight matrix, extends the per-node state
vectors $\lambda, u, r$ by one entry, and (if a downstream layer exists)
extends that layer's incoming dimension by one column. The new node is
initialized fully plastic ($\lambda_{\text{new}} = 0$,
$u_{\text{new}} = 0$, $r_{\text{new}} = 1$) so it can adapt freely while
existing nodes remain protected by EWC. The new row's incoming weights are
initialized along the top principal direction of the residual signal
$D = f(X_a) - f(X_b)$ at the layer below --- the direction of unmet
representational variance, computed by SVD over a probe batch.

\paragraph{Cellular pruning (\texttt{prune\_node}).}
Removes a node's row, contracts the state vectors, and drops the
corresponding input column on the next layer. By default the pruned node's
outgoing column is \emph{redistributed} to its weight-cosine-nearest peer
on the next layer before removal, preserving approximate input--output
behavior of the layer.

\paragraph{Growth trigger.}
A node is grown only when three conditions are jointly satisfied over a
window of $W$ steps:
\begin{enumerate}[leftmargin=2em]
  \item \emph{Loss plateau.} The loss has not improved by more than
        $\varepsilon_{\text{loss}}$ relative to the prior window.
  \item \emph{Rank saturation.} The effective rank
        \begin{equation*}
          \reff = \exp\!\left(-\sum_k p_k \log p_k\right),
          \qquad p_k = \sigma_k \big/ \textstyle\sum_k \sigma_k,
        \end{equation*}
        computed from the normalized singular values of the latent
        activation matrix, has approached the latent dimension
        ($\text{latent\_dim} - \reff < \varepsilon_{\text{rank}}$).
  \item \emph{Gradient stability.} The median gradient norm is bounded
        $g_{\min} \leq |\tilde g| \leq g_{\max}$ --- the optimizer is not
        exploding or vanishing.
\end{enumerate}
The conjunction of plateau-and-saturation-and-stability ensures that the
network grows only when it has \emph{actually run out of capacity for the
current task}, not merely when training is slow.

\paragraph{Resource ceilings.}
A pre-flight check rejects a proposed division if it would exceed a
per-curriculum byte budget \textit{cap\_bytes}. The cap can be lifted
between phases (we do this in the extension experiment of
\S\ref{sec:ship-wake-extend} and \S4.5). This implements \emph{resource
awareness} --- the network knows when not to grow even if growth would
help. Biologically analogous to metabolic limits on adult neurogenesis.

\paragraph{Epigenetic analogue.}
Growth corresponds to adult neurogenesis (which remains active in the
dentate gyrus and a few other regions throughout life); the
rank-saturation trigger reflects the biological observation that new
neurons are recruited preferentially when existing circuits saturate.
Pruning corresponds to developmental and ongoing synaptic pruning, with
cosine-nearest-peer redistribution as a soft analog of compensatory
hypertrophy.

\subsection{The Frustration Multiplier}
\label{sec:frustration}

We introduce a per-task plateau counter that scales the task loss when the
optimizer gets stuck. For each task, accumulate the loss in windows of
length $W_f$. At each window boundary, compare this window's mean loss to
the prior window's. If improvement is below $\varepsilon_f$, increment a
stuck counter $s$. The multiplier is
\begin{equation}
  m = \min\!\bigl(M_{\max},\; 1 + g \cdot \max(0, s - \tau + 1)\bigr)
\end{equation}
with hinge threshold $\tau$, gain $g$, and ceiling $M_{\max}$. The
multiplier is applied \emph{only to the task loss}, not the EWC penalty:
\begin{equation}
  L = m \cdot \Ltask + \beta \cdot \Lewc.
\end{equation}
This is mechanically equivalent to focal-loss / hard-example-mining:
amplify the gradient signal on examples the optimizer isn't progressing
on, without amplifying the regularizer.

\paragraph{Epigenetic analogue.}
The biological correlate is the hypothalamic--pituitary--adrenal (HPA)
axis stress response and its downstream effects on synaptic plasticity.
Sustained behavioral stress (failure to make progress on a task) elevates
glucocorticoids, which modulate gene expression at plasticity-related loci
through epigenetic mechanisms --- DNA-methylation changes at NR3C1 and
BDNF, miRNA-mediated post-transcriptional regulation of plasticity
proteins. Most epigenetic modifications require a stress signal to engage;
the frustration multiplier is the architecture's stress proxy.

\subsection{The Dreaming Phase}
\label{sec:dreaming-phase}

Between training tasks we run an offline \emph{dreaming block}.
Frustration is disabled inside the dream block (no environmental stress
during sleep); EWC remains active so that consolidation pulls weights
toward their anchored state. The dream block runs four sequential stages:
\begin{equation*}
  \text{replay} \;\to\; \text{compression (optional)} \;\to\;
  \text{purge} \;\to\; \text{archive}.
\end{equation*}
Replay is the load-bearing mechanism in the working configuration; the
compression-action design space is documented as ablation arms in
\S\ref{sec:compression-actions}.

\subsubsection{Replay}
\label{sec:replay}

Replay rehearses past tasks under EWC. We support two replay sources:
\begin{itemize}[leftmargin=1.5em]
  \item \emph{Episodic exemplars (hippocampal buffer).} When real-data
        storage is permitted we maintain a buffer of $K$ samples per past
        task, sampled uniformly at task end. During replay the buffer is
        shuffled with the current task's batch (1{:}1 by default) and the
        cross-entropy loss is evaluated on both. $K \in \{1, 10, 20, 50\}$
        is the storage axis we sweep in \S4.3. $K = 0$ routes replay
        through the storage-free pseudo-rehearsal pathway of
        \S\ref{sec:manifold-replay}.
  \item \emph{Manifold pseudo-rehearsal (storage-free).} Per-class
        statistics of L0 activations are stored at task end and used to
        synthesize replay samples; see \S\ref{sec:manifold-replay} for the
        full mechanism. This is the headline rehearsal channel.
\end{itemize}
Conceptually analogous to sharp-wave ripples during slow-wave sleep, which
replay recent and remote memories at compressed timescales.

\subsubsection{Redundancy Detection (Activation Cosine)}
\label{sec:redundancy}

When dream-time compression is enabled, we score each pair of nodes
$(i, j)$ within a layer for \emph{functional} redundancy via
centered-activation cosine over a probe batch $a_i, a_j \in \R^{B}$:
\begin{equation}
  \rho_{ij} = \frac{\langle a_i - \bar a_i,\; a_j - \bar a_j \rangle}
                   {\|a_i - \bar a_i\| \cdot \|a_j - \bar a_j\|}.
\end{equation}
Pairs with $\rho_{ij} \geq \tau_{\ac}$ are candidates for consolidation.
We use the activation signal (not weight-cosine) because it captures
functional similarity across the data distribution rather than identical
input-direction tuning. An early diagnostic confirmed that hidden-layer
weight-cosine ($0.6$--$0.8$) is high but not functional, while
latent-space weight-cosine is near zero by construction of the growth rule
(\S\ref{sec:structural-plasticity}) --- so the weight-cosine signal is
misleading. Layer-1 activation cosine reaches $0.85$--$0.92$ in practice,
well above the $0.59$--$0.81$ weight-cosine ceiling.

\subsubsection{The sRNA Cap}
\label{sec:srna-cap}

Each compression pass is bounded by a per-layer event cap $N_{\max}$. We
observe empirically (\S4) that $\sigma$-distance from baselines degrades
monotonically with event count, motivating small caps; $N_{\max} = 1$
allows replay to absorb each consolidation event before the next event
drifts on top.

\paragraph{Epigenetic analogue.}
The cap is a direct analog of resource-limited small-RNA pools that gate
sleep-cycle synaptic homeostasis. miRNAs involved in synaptic remodeling
are produced and consumed in stoichiometric amounts during a single sleep
cycle; the cellular machinery cannot perform arbitrarily many
consolidation events per cycle even when many candidates are present. The
``fewer events $=$ better'' empirical pattern reproduces the biological
observation that few consolidation events per cycle are protective.

\subsubsection{Compression-Action Design Space (Ablation Arms)}
\label{sec:compression-actions}

Three substrate-respecting mechanisms for resolving a detected redundant
pair $(i, j)$ were tested as dream-time compression actions. None of them
$\sigma$-confidently beats the no-compression baseline at 6 seeds; we
document them here as the ablation arms reported in \S4 and to motivate
the substrate-preserving emphasis of \S\ref{sec:archive}.

\begin{enumerate}[leftmargin=2em]
  \item \emph{Synaptic downscale} (substrate-preserving). The peer (kept
        side) absorbs the victim's outgoing column on the next layer; the
        victim's outgoing is zeroed; the victim's row at layer $L$ is
        \emph{not} touched. Architecture and parameter count are
        unchanged. Fisher entries on the victim's outgoing edge are reset
        so the dormant unit is available for re-recruitment in future
        tasks.
  \item \emph{Routing starvation} (asymmetric ramp). Each event multiplies
        the victim's routing-scale entry by $\alpha < 1$; below floor
        $\eta$ the scale latches to $0$ permanently. Bias is untouched, so
        the victim continues producing a constant downstream signal until
        downstream layers learn to absorb the constant via gradient
        descent --- the unit ``dies slowly'' rather than instantly.
        Reversible: each \texttt{compress()} pass also runs a regrow step
        on non-victim, non-latched units, multiplying their scale by
        $1/\alpha$ (capped at $1$). The kept side (``primary'') is
        selected by older \texttt{task\_of\_origin}, with larger
        outgoing-norm as tiebreaker.
  \item \emph{Apoptosis spike} (full-latch handler). When a
        routing-starvation event causes scale to cross $\eta$, two coupled
        mechanisms fire: (a) \emph{redistribution} --- the dead cell's
        outgoing column is transferred uniformly across all surviving
        non-latched peers, then zeroed; (b) \emph{spike} --- the surviving
        peers' apoptosis-pulse $p_k$ is raised to
        $p_{\text{init}} \in (0, 1]$. The pulse decays multiplicatively by
        $\delta$ per dream block. Through the $\lambda^{\text{eff}}$
        modification of \S\ref{sec:ewc-substrate}, neighbors of a fresh
        death train with reduced EWC stiffness for several dream cycles,
        allowing them to absorb the dead cell's role.
\end{enumerate}

We report engagement statistics for these mechanisms in \S4 as evidence
that the sRNA cap (\S\ref{sec:srna-cap}) is the load-bearing decision: at
the activation-cosine threshold $\tau_{\ac} = 0.92$, half the seeds fire
zero consolidation events --- a bimodal engagement pattern that
across-seed averaging conceals.

\subsubsection{Purge}
\label{sec:purge}

After compression, we drop nodes whose utility $u_i$ falls below a
threshold $\tau_u$. Reuses the structural pruning machinery of
\S\ref{sec:structural-plasticity}, including cross-layer fan-in cleanup
and cosine-nearest-peer redistribution. Distinct from the in-training
pruner (\S\ref{sec:structural-plasticity}): this is a sleep-time
structural sweep, not an event-clock check during waking training.

\subsubsection{Pseudo-code of a Dream Block}
\label{sec:dream-pseudocode}

\begin{verbatim}
for each task t in curriculum:
    train_one_task(t)                         # waking
    consolidate_task(t)                       # Fisher update + anchor
    if dreaming_enabled:
        decay apoptosis pulses by delta
        replay(buffer = hippo or manifold,
               steps = K_r,
               ewc_active = True)
        compress(signal = activation,
                 threshold = tau_ac,
                 action = ACTION,
                 max_events_per_layer = N_max)
        purge(threshold = tau_u)
        archive_block(...)                    # 3.7
\end{verbatim}

\subsection{Manifold Replay (Storage-Free Pseudo-Rehearsal)}
\label{sec:manifold-replay}

The headline rehearsal mechanism is \emph{manifold replay}: a
trioron-native pseudo-rehearsal scheme that summarizes each past class as
a small set of statistics in the network's first-hidden-layer (L0)
representation rather than storing exemplars.

\paragraph{Storage.}
When a task $t$ ends we collect L0 activations for each class $c$ that
appeared in the task and store the per-class mean
$\mu_c \in \R^{L_0}$ and per-dimension standard deviation
$\sigma_c \in \R^{L_0}$ (a diagonal-Gaussian summary). Storage per class is
$2 \cdot L_0 \cdot 4 = 1024$ bytes at $L_0 = 128$. A 30-class run costs
$\approx\!30$~KB total --- independent of the dataset size and of the
number of samples per task.

\paragraph{Replay.}
During the dream block we draw pseudo-samples for each past class $c$:
\begin{equation}
  \tilde z = \mu_c + \sigma_c \odot \varepsilon,
  \qquad \varepsilon \sim \mathcal{N}(0, I),
\end{equation}
then forward them through the \emph{upper} layers via
\texttt{forward\_from\_layer(z\_tilde, start = 1)} to obtain logits.
Cross-entropy loss on $(\tilde z, c)$ is added to the total loss. Replay
is performed at the same EWC strength as waking training.

\paragraph{Why per-class statistics work where logit-inversion fails.}
Three single-seed inversion variants --- DeepInversion-style code
synthesis (one-shot, refresh-all, and diversity-regularized) --- all
peaked at $0.27$ full-softmax accuracy on chained-15. Logit inversion
locates class-$c$ attractor \emph{peaks}, not class-$c$ manifold
\emph{samples}; the resulting points concentrate at adversarial extrema
that do not reflect actual class structure. This is the same pathology
that defeated earlier engram-replay variants (gradient-ascent on real-mean
seeds; L1-feature distillation; chain combinations of engram with
differential L1 storage). The $(\mu, \sigma)$ summary lives on the real L0
manifold by construction; inversion attractors do not. An ablation
removing per-class $\sigma$ ($\mu$-only replay) loses $\approx\!0.07$
absolute on full-softmax at 30 classes (single-seed), confirming that
per-class noise scale is load-bearing --- roughly $12\%$ of manifold
replay's advantage.

\paragraph{Why L0 is frozen.}
Manifold replay assumes that L0's encoding is approximately stationary
across the curriculum, so that statistics collected after task $t$ remain
valid at task $t + k$. We enforce this by freezing L0 at initialization to
a fixed random Gaussian projection (weights drawn from
$\mathcal{N}(0, 2/d)$ at init and never updated). A separation-probe
diagnostic on the chained-15 curriculum gives a class-prototype
discriminability ratio of $1.063$ in the frozen L0 space --- modest but
non-trivial, with Fashion-MNIST contributing the lowest separability
($1.04$) as expected from its irreducible within-domain confusions. The
ratio rules out a single-prototype scheme (L0-space class means alone are
not separable enough) and motivates the stored-$(\mu, \sigma)$ sketch as a
manifold summary rather than a class-prototype classifier.

\paragraph{Epigenetic analogue.}
Per-class $(\mu, \sigma)$ is a compact ``engram fingerprint'' --- an
offline sketch of the data manifold rather than a stored sample.
Functionally closer to consolidated memory traces (gist-level recall) than
to episodic exemplar storage (which the hippocampal-buffer baseline of
\S\ref{sec:replay} implements explicitly).

\subsection{The Dream Archive}
\label{sec:archive}

Long curricula accumulate parameters that should not keep moving --- rows
whose Fisher has plateaued and whose contribution is stable across recent
tasks. We expose these as a \emph{dream archive}: a two-phase mechanism
for locking and quantizing such rows.

\paragraph{Phase 1 --- mid-curriculum row archive.}
After each \texttt{consolidate\_task}, every row of every layer is scored
along three axes: (a) $\lambda$ percentile (top
\textit{archive\_lam\_top\_percentile} of the layer), (b)
gradient-magnitude ceiling (median $|g|$ over the recent window below
\textit{archive\_grad\_mag\_floor}), and (c) apoptosis-pulse bound
($p \leq$ \textit{archive\_pulse\_max} --- never archive a row recently
pulsed by an apoptosis event). A row that satisfies all three for
\textit{archive\_streak\_threshold} consecutive tasks is \emph{archived}:
a per-row gradient mask is set so subsequent backward passes leave the row
at its anchor value. The mask is enforced via
\texttt{mask\_archived\_grads\_all()} between \texttt{loss.backward()} and
\texttt{opt.step()} (sequenced after PackNet/HAT grad surgery on
competitor arms). A per-layer cap \textit{archive\_max\_per\_layer}
prevents any one layer from being fully archived in a single pass.

\paragraph{Phase 2 --- end-of-curriculum int8 quantization.}
At end of training we snap each archived row to int8 with a per-row
symmetric scale ($8$~bits/weight plus one fp32 scale per row). Active
(non-archived) weights are rounded to BF16 precision. The Phase 2 path is
implemented as a snapshot--quantize--re-evaluate--restore sequence in the
bench so we can report storage and accuracy numbers from the same trained
checkpoint; the deployment path applies the snap permanently.

\paragraph{Storage accounting.}
With BF16 as the deployment baseline, archive $+$ int8 reduces network
weight storage by $\approx\!40\%$ across all arms ($213.7 \to 127.5$~KB on
the headline \texttt{grown\_capped\_dream} arm; full table in \S4.4). We
report deployment storage in the BF16 baseline because that is the honest
device-side accounting; an FP32 baseline overstates the dream-archive win.

\paragraph{Why ternary fails on a frozen-L0 substrate.}
Ternary quantization ($2$~bits/weight) on the frozen-Gaussian L0
projection costs $0.31$ absolute full-softmax accuracy on the smoke test:
random-Gaussian weights have no concentration of magnitude near zero, so
$2$-bit quantization throws away most of the routing structure that L1
and the head learned to depend on. Int8 ($\approx 1/127 \approx 0.8\%$
per-weight noise) preserves the routing intact. Unlocking ternary would
require quantization-aware training during the dream phase, or restricting
ternary to layers whose weights are concentrated at zero (a learned,
pruned L1 + head) --- both deferred to future work.

\paragraph{Epigenetic analogue.}
Phase 1 corresponds to terminal cell differentiation --- a cell that has
settled into its mature role no longer participates in plasticity. Phase 2
corresponds to compression of consolidated memory traces into
low-bandwidth long-term storage, freeing high-bandwidth substrate for new
learning.

\subsection{Mixed-Precision Substrate}
\label{sec:mixed-precision}

For deployment on resource-constrained devices the architecture supports a
mixed-precision mode in which weight Parameters $(W, b)$ are stored in
narrow precision (BF16) while all consolidation buffers ---
$\anchorW, \anchorB, F_W, F_b, \lambda, u, r, p$ --- remain in FP32. The
forward pass auto-casts inputs to the weight dtype, preserving the
callable interface; the EWC penalty and Fisher accumulation upcast cleanly
across the boundary.

This separation is necessary: pure-narrow-precision training degrades by
$\approx\!40\%$ due to optimizer-state precision loss in Adam's running
moments, but inference latency and memory cost are dominated by the
weights, which can safely run in BF16. The intended deployment pattern is
\emph{train in FP32 during sleep cycles, deploy in BF16 for always-on
inference}.

\subsection{The Ship-Wake-Extend Loop}
\label{sec:ship-wake-extend}

Deployment is not a one-shot training run that yields a fixed artifact; it
is an iterative loop in which the model is shipped, woken on new data,
allowed to grow plastic substrate, and re-shipped. We instantiate this
loop with a \emph{consolidation-dream pass} that fires once at a chosen
task boundary:
\begin{enumerate}[leftmargin=2em]
  \item \texttt{consolidation\_dream\_pass} --- full-coverage replay over
        all past tasks with archive-aware gradient masking (locks
        honored), plus one additional \texttt{archive\_block} call to
        catch rows that just settled.
  \item Optional permanent int8 snap of the archived rows (no restore
        step, since this is the shipping checkpoint).
  \item Cap lift:
        $\textit{current\_cap\_bytes} \leftarrow \textit{extension\_cap\_bytes}$.
        The network is allowed to grow new plastic substrate against the
        new task budget.
  \item The training loop continues over the new tasks with shared
        internal state --- manifold buffer, EWC anchors, archived rows,
        frustration counters all persist across the boundary.
\end{enumerate}

The result is a smaller plastic envelope in deployment than in training,
with new growth allocated only against the residual. A device-conscience
deployment cannot revisit the original training data: it has only what it
has stored (the manifold buffer, the archived weights) and what arrives in
the new data stream. The consolidation-dream pass turns shipping into a
controlled handoff: the network ``sleeps long'' before going to market,
locking in the rows it will never need to move again, then wakes on new
data in the field.

\subsection{Experiment Design}
\label{sec:experiment-design}

\paragraph{Curriculum (chained-15).}
A 30-class class-incremental benchmark drawn from three datasets in
sequence:
\begin{itemize}[leftmargin=1.5em]
  \item 10 MNIST digit tasks (one new digit class per task);
  \item 10 Fashion-MNIST tasks (one new garment class per task);
  \item 10 EMNIST-letters tasks A..J (one new letter class per task).
\end{itemize}
Each task introduces exactly one new class. At evaluation time the network
is asked to classify across all classes seen so far (full-softmax) or
restricted to the same task or domain (task-aware, domain-aware). The
class-incremental framing forces the network to commit to absolute class
identities --- there is no task ID at inference --- and is the standard
hard regime for continual-learning evaluation.

\paragraph{Why class-incremental, not task-incremental.}
Task-incremental benchmarks supply the task ID at inference, which
collapses the problem to learning a per-task classifier.
Class-incremental does not, and is therefore the regime where
partition-style methods (PackNet, HAT) suffer their characteristic
full-softmax collapse. Reporting class-incremental from the start avoids
confusing the comparison.

\paragraph{Extension curriculum (chained-23).}
chained-15 above, plus 8 additional EMNIST-letters tasks K..R introduced
at the extension boundary; total 23 class-incremental tasks across 38
classes. Used only in the extension experiment of \S4.5.

\paragraph{Network configuration.}
\begin{itemize}[leftmargin=1.5em]
  \item Three layers: input $\to$ L0 $\to$ L1 $\to$ head.
  \item L0: $784 \to 128$, frozen-Gaussian random projection (weights
        drawn from $\mathcal{N}(0, 2/784)$ at init and never updated). See
        \S\ref{sec:manifold-replay} for why L0 is frozen.
  \item L1: $128 \to 32$ initial, grown adaptively up to
        $\textit{cap\_bytes} = 8\,\text{K}$ trainable parameters during
        chained-15, lifted to 16~K at the extension boundary.
  \item Head: $L_1\text{\_width} \to n_{\text{classes\_so\_far}}$, grown
        adaptively as new classes appear.
  \item Activations: ReLU on L1, identity on the head.
  \item Optimizer: Adam, $\text{lr} = 3 \times 10^{-3}$.
  \item Epochs: 4 per task. Loss: cross-entropy on the active label set.
  \item EWC strength: $\beta = 1000$ during waking and dream-replay
        (matched).
\end{itemize}

\paragraph{Trioron arms.}
\begin{itemize}[leftmargin=1.5em]
  \item \emph{fixed\_ewc\_small} --- no growth; L1 fixed at the
        chained-15 cap. Isolates the EWC + manifold-replay contribution at
        matched parameters.
  \item \emph{grown\_capped\_no\_dream} --- growth under EWC, no
        dream-time consolidation. Isolates the dream contribution.
  \item \emph{grown\_capped\_dream} --- growth + dream within cap. The
        default deployment arm.
  \item \emph{grown\_uncapped\_dream} --- growth past cap with manifold
        replay. The upper-bound headline arm.
\end{itemize}

\paragraph{Competitor families (5).}
We benchmark against one method from each major continual-learning
family:
\begin{itemize}[leftmargin=1.5em]
  \item \emph{Hippocampal buffer (exemplar rehearsal).} $K$ samples per
        past class stored. $K \in \{1, 10, 20, 50\}$.
  \item \emph{PackNet} (Mallya \& Lazebnik 2018). Iterative magnitude
        pruning with per-task masks. \emph{PackNet matched} matches our
        trainable-parameter budget; \emph{PackNet standard} uses the
        original recipe.
  \item \emph{HAT} (Serra et~al. 2018). Hard binary attention masks per
        task, gating each layer. \emph{Matched} and \emph{standard}
        variants as above.
  \item \emph{Online EWC} (Schwarz et~al. 2018). Fisher EMA across tasks
        with $\gamma = 0.95$; learning-rate and $\beta$ tuned via grid
        search on the no-dream loss.
  \item \emph{LwF + EWC} (Li \& Hoiem 2018; Kirkpatrick 2017).
        Logit-distillation toward the anchored network on old-class
        columns, in addition to the EWC penalty.
\end{itemize}

\paragraph{Metrics.}
\begin{itemize}[leftmargin=1.5em]
  \item \emph{Full-softmax accuracy} (full): mean across all classes seen
        so far of top-1 accuracy with the unrestricted softmax. The hard
        class-incremental metric.
  \item \emph{Domain-aware accuracy} (domain): top-1 accuracy restricted
        to classes that share a dataset domain (digits, fashion, letters)
        with the query. Measures preservation of coarse-grained
        discrimination.
  \item \emph{Task-aware accuracy} (task): top-1 accuracy restricted to
        the classes of the \emph{task} the query belongs to. The
        pairwise-discrimination metric; the right one for trioron's
        deployment role of context-gated personalization atop a frozen
        LLM.
  \item \emph{Forgetting} (forget): mean across past tasks of (final-task
        accuracy on that task) $-$ (accuracy at the time the task was
        first introduced); positive values indicate forgetting, negative
        values indicate \emph{negative forgetting} (accuracy improves over
        time after the task ended).
  \item \emph{$\sigma$-vs-baseline}: paired-by-seed standardized effect
        size across $N = 3$ random seeds, reported alongside absolute
        $\Delta$.
\end{itemize}

\paragraph{Storage accounting.}
We report deployment storage in two columns: BF16 baseline (the honest
device-side number, since BF16 is the inference dtype) and FP32 baseline
(the training-time number). The manifold buffer is counted at FP32
(per-class $\mu$ and $\sigma$, $2 \cdot L_0 \cdot 4 = 1024$ bytes/class).
All headline storage numbers in \S4 use the BF16 baseline.

\paragraph{Reproducibility.}
All benches run on commodity CPU. Each seed is fully deterministic given
the seed value, modulo a known torch nondeterminism in archive-trigger
timing (see \S4.4 for the v1/v2 reconciliation). Launcher scripts under
\texttt{experiments/} accept explicit seed lists; bench logs and CSVs from
every reported run are committed to the public repository.

\end{document}
